% ================================================================= 1. TÓM TẮT
\section{Tóm tắt}

Theo dõi tim thai không xâm lấn bằng \textbf{một} điện cực dán trên bụng mẹ là cấu hình rẻ nhất và
dễ đeo nhất cho giám sát tại nhà, đồng thời là cấu hình khó nhất về mặt xử lý tín hiệu. Sóng tim thai
bị sóng tim mẹ lớn gấp nhiều lần che khuất, hai nguồn trùng phổ tần số, và khi chỉ có một kênh thì
không dùng được các phương pháp tách nguồn mù đa kênh vốn cho kết quả tốt nhất.

Hai mươi năm nghiên cứu đã sinh ra hàng chục phương pháp với F1 công bố trải từ 77\% đến 99,7\%.
Nhóm đã đọc toàn văn 30 công trình và kết luận rằng \textbf{phần lớn các con số đó không so sánh
được với nhau}, vì chúng khác nhau ở cách tách chủ thể, dung sai ghép cặp, số kênh thực sự dùng, và
việc có tinh chỉnh trên tập đích hay không. Chi tiết ở \S\ref{sec:yvan}.

Trước khi viết đề cương này, nhóm đã chạy một loạt thí nghiệm tiền khả thi. Ba kết quả định hình
toàn bộ phần còn lại của tài liệu.

\begin{enumerate}[leftmargin=1.5em,itemsep=4pt]
\item \textbf{Front-end tín hiệu quan trọng hơn kiến trúc mạng rất nhiều.}
Chỉ đổi dải thông từ 1--45\,Hz sang 10--60\,Hz nâng macro F1 thêm \tot{11,00 điểm} trong cùng một
thí nghiệm có kiểm soát. Trong khi đó, khi giữ cố định ngữ cảnh thời gian, việc thay gradient boosting
bằng mạng học sâu chỉ đáng \xau{+0,41 điểm} với $p = 0{,}7012$, tức không có ý nghĩa thống kê.

\item \textbf{Biểu diễn topo không phù hợp cho bài toán dò đỉnh.}
Persistence Image kết hợp CNN 2D đứng \xau{cuối cùng trong 16 kiến trúc} được so ở cùng ngân sách
tham số (F1 27,42 so với 97,14 của CNN 1D cùng 27 nghìn tham số). Nhóm cũng chứng minh được bằng số
rằng đặc trưng $H_0$ của lọc dưới mức trên tín hiệu một chiều \textbf{trùng khớp 100\%} với độ nhô đỉnh
cổ điển, tức không phải một đặc trưng mới.

\item \textbf{Rào cản thật không phải độ chính xác trung bình mà là độ tin cậy.}
Mô hình của nhóm đạt macro F1 \tot{99,21} trên ADFECGDB khi tách theo bản ghi, nhưng khi chuyển sang
PhysioNet/CinC 2013 mà không tinh chỉnh thì chỉ còn \xau{77,34}. Quan trọng hơn con số trung bình là
hình dạng phân bố: bốn trên mười bản ghi đạt F1 tuyệt đối 100\%, ba bản khác sụp xuống dưới 50\%.
\textbf{Hệ thống không biết khi nào nó sai.}
\end{enumerate}

Từ ba quan sát đó, đề tài đề xuất \hethong{}: một bộ dò QRS thai đơn kênh \textbf{kèm chỉ số chất lượng
tín hiệu thai dựa trên đồng điều bền vững}, cho phép hệ thống \textbf{từ chối trả lời} khi tín hiệu
không đủ tin cậy, thay vì im lặng đưa ra một nhịp tim sai.

\begin{ghichu}[title={Ba đóng góp}]
\begin{description}[leftmargin=2.2em,style=nextline,itemsep=3pt]
\item[C1 --- Đối chuẩn không rò rỉ.] So sánh năm họ biểu diễn tín hiệu dưới cùng một giao thức,
cùng bộ phân loại, cùng ngân sách huấn luyện, trên ba bộ dữ liệu, với quy tắc chọn kênh
\textbf{mù nhãn} có thể đăng ký trước.
\item[C2 --- Phân rã đóng góp ba tầng.] Trả lời định lượng câu hỏi: trong toàn bộ pipeline, phần nào
thực sự tạo ra hiệu năng? Đây là xương sống khoa học của đề tài.
\item[C3 --- Chỉ số chất lượng tín hiệu thai dựa trên đồng điều bền vững, có quyền từ chối trả lời.]
Đây là ứng dụng đầu tiên của phân tích dữ liệu topo cho tín hiệu tim thai mà nhóm tìm được.
\end{description}
\end{ghichu}

\clearpage

% ================================================================= 2. ĐẶT VẤN ĐỀ
\section{Đặt vấn đề}

\subsection{Bối cảnh lâm sàng}

Nhịp tim thai và hình thái phức bộ QRS thai là cửa sổ chính để phát hiện suy thai, thiếu oxy và một số
bất thường tim bẩm sinh. Hai phương pháp đang dùng trên lâm sàng đều có hạn chế rõ ràng.

\begin{table}[H]\centering\small
\caption{Hai phương pháp theo dõi tim thai hiện hành và giới hạn của chúng.}
\label{tab:lamsang}
\begin{tabular}{@{}L{4.0cm} L{5.0cm} L{5.4cm}@{}}
\toprule
\textbf{Phương pháp} & \textbf{Ưu điểm} & \textbf{Giới hạn} \\
\midrule
Tim thai đồ / Doppler (CTG) & Không xâm lấn, phổ biến & Chỉ cho nhịp, không cho hình thái. Cần kỹ thuật
viên đặt đầu dò. Không dùng liên tục tại nhà được. \\
Điện cực xoắn da đầu thai & Chuẩn vàng, tín hiệu sạch & Xâm lấn. Chỉ dùng được sau khi đã chuyển dạ
và vỡ ối. Không dùng cho theo dõi thai kỳ. \\
\bottomrule
\end{tabular}
\end{table}

Điện tim thai không xâm lấn từ điện cực ổ bụng cho phép theo dõi liên tục, tại nhà, chi phí thấp, và
cho cả hình thái sóng chứ không chỉ nhịp. Cấu hình \textbf{đơn kênh} là rẻ nhất và dễ tích hợp vào đai
đeo nhất, nhưng cũng loại bỏ luôn khả năng dùng tách nguồn mù đa kênh --- vốn là họ phương pháp đạt độ
chính xác cao nhất trong đối chuẩn tổng hợp của Andreotti và cộng sự.

\subsection{Vì sao đơn kênh khó}

Bảng dưới đây là các con số nhóm \textbf{tự đo được} trên ADFECGDB, không phải trích từ y văn.

\begin{table}[H]\centering\small
\caption{Các khó khăn của cấu hình đơn kênh, kèm hệ quả định lượng nhóm đo được trên ADFECGDB.}
\label{tab:khokhan}
\begin{tabular}{@{}L{6.2cm} L{8.2cm}@{}}
\toprule
\textbf{Thách thức} & \textbf{Hệ quả định lượng} \\
\midrule
Biên độ sóng thai nhỏ hơn sóng mẹ nhiều lần & Tỉ số RMS giữa vùng QRS thai và nền chỉ \textbf{1,24}
ở dải 1--45\,Hz; lên \textbf{1,64} ở dải 15--60\,Hz \\
Trùng phổ tần số trong khoảng 0,5--40\,Hz & Lọc tuyến tính không tách được nếu không làm méo sóng thai \\
Nhịp mẹ và nhịp thai trùng pha & \textbf{17--19\%} nhịp thai nằm trong $\pm$60\,ms của một nhịp mẹ \\
Không có đa dạng không gian & Không dùng được tách nguồn mù đa kênh (ICA, PCA, $\pi$CA) \\
Chất lượng kênh biến thiên cực lớn & F1 theo từng kênh dao động từ \textbf{85,03} đến \textbf{100,00}
trên cùng một bản ghi \\
Dữ liệu có nhãn rất ít & ADFECGDB chỉ có 5 sản phụ; chia ngẫu nhiên là rò rỉ dữ liệu \\
\bottomrule
\end{tabular}
\end{table}

\subsection{Ba khoảng trống trong y văn}
\label{sec:khoangtrong}

Nhóm đã tải và đọc \textbf{toàn văn 30 công trình} (danh mục đầy đủ ở \S\ref{sec:yvan} và Phụ lục).
Ba khoảng trống được xác nhận.

\begin{description}[leftmargin=2.6em,style=nextline,itemsep=6pt]

\item[G1 --- Không có so sánh công bằng giữa các họ biểu diễn tín hiệu.]
Mỗi công trình dùng một biểu diễn riêng (tín hiệu thô, ma trận trễ, phổ đồ thời gian--tần số, đặc trưng
topo) trên dữ liệu riêng với giao thức riêng. Chưa có công trình nào giữ cố định đồng thời bộ phân loại,
dữ liệu, giao thức và ngân sách huấn luyện để so các biểu diễn. Hệ quả là cộng đồng không biết cải tiến
đến từ biểu diễn hay đơn giản là từ tiền xử lý tốt hơn.

\item[G2 --- Không có đánh giá độ tin cậy ở mức bản ghi.]
Mọi công trình đều báo cáo một con số F1 trung bình. Không công trình nào trả lời câu hỏi lâm sàng
thật sự: \textit{``với bản ghi cụ thể đang nằm trước mặt, đầu ra này có đáng tin không?''} Đo đạc của
nhóm cho thấy phân bố hiệu năng là \textbf{lưỡng cực}, nên chính con số trung bình đang che giấu thông
tin quan trọng nhất.

\item[G3 --- Chưa có công trình nào áp dụng đồng điều bền vững cho tín hiệu tim thai.]
Nhóm truy vấn có cấu trúc trên Europe PMC REST, Semantic Scholar, OpenAlex và arXiv API với tổ hợp
\texttt{("persistent homology" OR "topological data analysis" OR "persistence image") AND (fetal OR
foetal OR "abdominal ECG" OR cardiotocography)}, tìm bổ sung bằng tiếng Trung, tiếng Việt và tiếng
Tây Ban Nha. Không có kết quả nào là công trình về tín hiệu thai. Kết quả gần nhất là biến thiên nhịp
tim nhi khoa, mà chính tác giả nêu \textit{``mở rộng khung này sang fetal HRV''} là hướng tương lai.

\end{description}

\subsection{Câu hỏi nghiên cứu và giả thuyết}

\begin{table}[H]\centering\small
\caption{Sáu câu hỏi nghiên cứu và thí nghiệm trả lời tương ứng.}
\label{tab:rq}
\begin{tabular}{@{}C{1.1cm} L{10.6cm} L{2.5cm}@{}}
\toprule
\textbf{Mã} & \textbf{Câu hỏi} & \textbf{Thí nghiệm} \\
\midrule
RQ1 & Dưới một giao thức không rò rỉ, họ biểu diễn tín hiệu nào cho hiệu năng dò fQRS đơn kênh tốt nhất
khi giữ cố định bộ phân loại và ngân sách huấn luyện? & E1 \\
RQ2 & Trong toàn bộ pipeline, tầng nào đóng góp nhiều nhất: front-end tín hiệu, độ dài ngữ cảnh, họ
kiến trúc, hay hậu xử lý? & E2 \\
RQ3 & Chỉ số chất lượng dựa trên đồng điều bền vững có dự báo được khi nào bộ dò sẽ sai, và có tốt hơn
các chỉ số cổ điển (SampEn, bSQI, kurtosis, entropy phổ) không? & E3 \\
RQ4 & Cơ chế từ chối trả lời nâng độ tin cậy lên bao nhiêu, đánh đổi bằng bao nhiêu phần trăm dữ liệu
bị loại? & E4 \\
RQ5 & Hiệu năng suy giảm thế nào khi chuyển miền và khi tỉ số tín hiệu trên nhiễu giảm? & E5, E6 \\
RQ6 & Tiền huấn luyện trên dữ liệu tổng hợp rồi tinh chỉnh trên dữ liệu thật có cải thiện khả năng
tổng quát hoá không? & E7 \\
\bottomrule
\end{tabular}
\end{table}

\begin{table}[H]\centering\small
\caption{Giả thuyết kiểm định được và tiêu chí chấp nhận.}
\label{tab:h}
\begin{tabular}{@{}C{1.1cm} L{9.0cm} L{4.2cm}@{}}
\toprule
\textbf{Mã} & \textbf{Giả thuyết} & \textbf{Tiêu chí chấp nhận} \\
\midrule
H1 & Bộ dò \hethong{} đạt F1 $\geq$ 95\% trên ADFECGDB khi tách theo bản ghi, với quy tắc chọn kênh
mù nhãn & Wilcoxon ghép cặp so với nền trừ mẫu kết hợp PCA, $p < 0{,}05$ \\
H2 & Front-end tín hiệu đóng góp nhiều hơn họ kiến trúc ít nhất 5 lần về điểm F1 & So sánh hai hiệu ứng
chính, kèm khoảng tin cậy bootstrap \\
H3 & fSQI topo có tương quan Spearman $|\rho| \geq 0{,}6$ với F1 thực theo bản ghi, và không kém
SampEn & Bootstrap CI cho hiệu số $\rho$ \\
H4 & Từ chối 20\% dữ liệu có fSQI thấp nhất nâng F1 trên phần còn lại thêm $\geq$ 8 điểm & Đường cong
rủi ro--độ phủ, so với cổng từ chối ngẫu nhiên \\
H5 & Tiền huấn luyện FECGSYNDB nâng F1 xuyên bộ dữ liệu thêm $\geq$ 5 điểm & Wilcoxon ghép cặp theo
bản ghi \\
H0 & \textbf{Giả thuyết phủ định vẫn có giá trị}: nếu H3--H5 không đạt, thì đối chuẩn cộng phân rã
đóng góp cộng kết quả phủ định có kiểm định vẫn là một bài báo hoàn chỉnh & --- \\
\bottomrule
\end{tabular}
\end{table}

\clearpage
